\noindent\makebox[0pt][l]{\fontsize{13}{16}\selectfont 418}\hfill{\small\textbf{NACHLASS.}}\hfill\null\par\medskip

\begin{center}
Historischer Bericht\\
über die von dem Hofrath \textsc{Gauss} theils ausgeführten, theils geleiteten\\
Messungen im Königreich Hannover.
\end{center}

\par\noindent\hspace*{\parindent}Die verschiedenen Messungsarbeiten, von welchen ich hier einen kurzen historischen Bericht abzustatten habe, sind zwar unter einander enge verknüpft, haben aber ungleiche Zwecke und ungleichen Charakter. Es wird am übersichtlichsten sein, sie nach der Zeitfolge der Aufträge zu ordnen, denen gemäss ich sie auf mich genommen habe.

\par\medskip
\begin{center}
I.\quad Die hannoversche Gradmessung.
\end{center}

\par\noindent\hspace*{\parindent}Durch Rescript vom 30. Junius 1820 wurde ich beauftragt, eine Gradmessung durch das Königreich Hannover auszuführen, als eine Erweiterung oder Fortsetzung einer nicht lange vorher angefangenen ähnlichen Arbeit in den dänischen Staaten. Man versteht unter jener Benennung diejenigen, theils astronomischen, theils trigonometrischen Operationen, wodurch die Grösse eines Meridiangrades in einem bekannten Längenmaasse (Fuss, Toise, etc.) bestimmt wird. Man wählt zu dem Ende zwei hinlänglich von einander entfernte Punkte in einerlei Meridian, bestimmt die Länge des zwischen ihnen enthaltenen Bogens in Fussen, etc. durch ein zwischen ihnen geführtes Dreiecksnetz, und die Anzahl von Graden, Minuten und Secunden, welche demselben Bogen entsprechen, durch astronomische an den Endpunkten angestellte Beobachtungen, woraus man dann auf die Länge Eines Grades zurückschliesst.

\par\noindent\hspace*{\parindent}Es erhellt hieraus, dass eine Gradmessung, \textls[120]{als solche}, nur die rein wissenschaftliche Tendenz hat, einen Beitrag zur mathematischen Kenntniss der Verhältnisse des Erdsphäroids zu geben, die nach den Anforderungen des Jahrhunderts eine viel grössere Schärfe und eine viel grössere Zahl von Gradmessungen nothwendig macht, als womit man früher sich begnügen musste. Es erhellt ferner hieraus, dass eine Gradmessung aus zweien sehr heterogenen Theilen besteht, einem trigonometrischen und einem astronomischen. Ohne hier in ein weiteres Detail einzugehen, will ich nur bemerken, dass ich den (25 Hauptdreiecke von Göttingen bis Hamburg umfassenden) trigonometrischen Theil in den Jahren 1821, 1822, 1823, und den astronomischen im Jahr 1827

\par\clearpage
\noindent\null\hfill{\small\textbf{ZUR HANNOVERSCHEN TRIANGULATION.}}\hfill\makebox[0pt][r]{\fontsize{13}{16}\selectfont 419}\par\medskip

\noindent ganz vollendet habe, und dass der letztere sowie die Endresultate in einem 1828 von mir herausgegebenen Werke »Bestimmung des Breitenunterschiedes zwischen den Sternwarten von Göttingen und Altona durch Beobachtungen am \textsc{Ramsden}schen Zenithsector« bekannt gemacht sind. Die Sternwarten von Göttingen und Altona, die durch ein einziges Spiel des Zufalls genau in einerlei Meridian liegen, sind nemlich die Endpunkte des gemessenen Meridianbogens, dessen ganze Krümmung etwas über zwei Grad beträgt.

\begin{center}
II.\quad Trigonometrische Verbindung der Gradmessungsdreiecke mit den Drei-\\
ecken der königlich niederländischen Vermessungen.
\end{center}

\par\noindent\hspace*{\parindent}Der nächste wissenschaftliche Zweck dieser mir durch Rescript vom 8. März 1824 aufgetragenen Erweiterung der trigonometrischen Arbeiten war, die hannoversche Gradmessung mit der französischen und englischen in Verbindung zu bringen. Diese letztern sind bekanntlich unter sich verbunden, und an die französischen Dreiecke schliessen sich die mit vieler Sorgfalt gemessenen Dreiecke in den Niederlanden an, welche in der Zeit, wo Ostfriesland mit Holland vereinigt war, bis an die östliche Grenze dieses Fürstenthums, weiter südlich hingegen bis nach Bentheim fortgeführt waren. Ich hatte also die Wahl unter zwei Wegen, auf denen sich der vorgesetzte Zweck erreichen liess, entweder nemlich von den nördlichsten Dreiecken der Gradmessung über Bremen nach Jever, oder von den mittlern durch Westphalen nach Bentheim. Der erstere Weg wurde ausser andern Gründen auch deswegen vorgezogen, weil dadurch zugleich eine Verbindung mit der Nordsee, und damit die Bestimmung der absoluten Höhen sämmtlicher Dreieckspunkte über dem Meeresspiegel, erreicht werden konnte. Diese Dreiecksmessung wurde in den Jahren 1824 und 1825 von mir ausgeführt, und damit die Zahl sämmtlicher Hauptdreiecke auf 38[\textsuperscript{*})] gebracht.

\par\medskip
\par\noindent\hspace*{\parindent}Über diese trigonometrischen Messungen von 1821---1825 habe ich noch einige Bemerkungen beizufügen, da sie sich von ähnlichen Arbeiten, wie früher in andern Ländern ausgeführt sind, in mehrern Beziehungen unterscheiden.

\par\smallskip\noindent\rule{.24\linewidth}{.4pt}\par\smallskip
{\footnotesize\noindent\hspace*{\parindent}[\textsuperscript{*}) Die Anzahl der unabhängigen Dreiecke der Gradmessung zwischen Göttingen und Hamburg beträgt 21, die Anzahl sämmtlicher unabhängigen Dreiecke der Gradmessung und ihrer Fortsetzung nach Jever (ohne das Dreieck Hohehagen-Brocken-Inselsberg) 42; vergl. S. 297.]\par}

\par\smallskip\noindent\hfill{\footnotesize 53\textsuperscript{*}}\hspace*{.12\linewidth}\par

\par\clearpage
\noindent\makebox[0pt][l]{\fontsize{13}{16}\selectfont 420}\hfill{\small\textbf{NACHLASS.}}\hfill\null\par\medskip

\par\noindent\hspace*{\parindent}1) Durch die Anwendung der von mir zuerst eingeführten \textls[120]{Heliotrope} wurden besondere Signalthürme entbehrlich, Dreiecke von einer früher impraktikabeln Grösse möglich, und eine ohne jenes Hülfsmittel nicht zu erlangende Schärfe der Messungen erreichbar.

\par\noindent\hspace*{\parindent}2) Bei der Ausführung der Messungen habe ich mich nicht auf das zu dem unmittelbaren Zwecke erforderliche (was, wie schon bemerkt ist, zunächst nur wissenschaftliche Tendenz hatte) eingeschränkt, sondern jene zugleich für die Landesgeographie so fruchtbar zu machen gesucht, wie nur, ohne dem nächsten Zwecke Abbruch zu thun, geschehen konnte. Es sind daher auch die meisten im Bereich der Dreieckspunkte liegenden Thürme sehr scharf festgelegt, nicht bloss von den in den Dreieckszug fallenden Städten, wie Göttingen, Hildesheim, Wolfenbüttel, Braunschweig, Hannover, Celle, Lüneburg, Harburg, Hamburg, Buxtehude, Stade, Verden, Bremen, Oldenburg, Varel, Jever, sondern auch von vielen hundert kleinern Ortschaften, wie die meinen Berichten über die Arbeiten von jedem Jahre beigefügten Übersichtskarten zeigen. Dadurch sind also feste Anhaltspunkte und Grundlagen für alle später in den betreffenden Landstrichen vorzunehmenden Detailaufnahmen gewonnen, und diese Bestimmungen erhalten dadurch einen bleibenden Werth.

\par\noindent\hspace*{\parindent}3) Endlich sind diese Resultate durch eine mir eigenthümliche Behandlungsweise in eine solche Form gebracht, die für die eben ausgesprochene weitere Benutzung wesentliche Vortheile darbietet.

\begin{center}
III.\quad Trigonometrische Vermessung der Landestheile,\\
welche von den Messungen 1821---1825 nicht berührt waren.
\end{center}

\par\noindent\hspace*{\parindent}Die in den Jahren 1821---1825 ausgeführten Arbeiten enthielten (wenn gleich nicht zunächst für diesen Zweck bestimmt) eine wirkliche trigonometrische Vermessung eines sehr beträchtlichen Theils des Königreichs Hannover: die allgemein anerkannten Vortheile, welche eine genaue trigonometrische Landesvermessung gewährt, liessen es als wünschenswerth erscheinen, dass hiebei nicht stehen geblieben würde. Durch Rescript vom 28. April 1828[\textsuperscript{*})] wurde mir der Auftrag ertheilt, die weitere Erstreckung der trigonometrischen Vermessung über alle durch die frühern Arbeiten noch nicht berührten Landestheile zu leiten.

\par\smallskip\noindent\rule{.24\linewidth}{.4pt}\par\smallskip
{\footnotesize\noindent\hspace*{\parindent}[\textsuperscript{*}) Das im Gauss-Archiv befindliche Rescript ist vom 14. April datirt.]\par}

\par\clearpage
\noindent\null\hfill{\small\textbf{ZUR HANNOVERSCHEN TRIANGULATION.}}\hfill\makebox[0pt][r]{\fontsize{13}{16}\selectfont 421}\par\medskip

\par\noindent\hspace*{\parindent}Zu gleicher Zeit wurde eine Detailaufnahme angeordnet, die alle Landestheile umfassen sollte, welche in der in den achtziger Jahren des vorigen Jahrhunderts ausgeführten Messtischaufnahme des vormaligen Kurfürstenthums Hannover noch nicht begriffen gewesen waren. Mit der speciellen Leitung dieser Detailaufnahme wurde der Oberst \textsc{Prott} beauftragt; diese Arbeit steht aber mit der trigonometrischen Vermessung insofern in unmittelbarer Verbindung, als die Resultate der letztern zur Grundlage und zu festen Anhaltspunkten für die Messtischarbeiten dienen. Welche Leichtigkeit und Sicherheit die Messtischarbeiten auf diese Weise gewinnen, hatte sich schon ein Jahr zuvor bei einigen vorläufigen Detailaufnahmen im Hildesheimschen bewährt, obgleich damals die Ausführung durch Officiere geschah, denen früher diese Methode fremd gewesen war.

\par\noindent\hspace*{\parindent}Was ich nun hier über die unter meine Leitung gestellte trigonometrische Messung zu sagen habe, betrifft theils das dabei thätig gewesene Personal, theils den Umfang der bisher vollendeten, theils endlich die nun noch rückständigen Messungen.

\begin{center}
A.\quad Das Personal.
\end{center}

\par\noindent\hspace*{\parindent}Zu einer trigonometrischen Messung sind zweierlei ganz verschiedenartige Arbeiten erforderlich, die Ausführung der Messungen an den betreffenden Plätzen im Felde, und ihre Verarbeitung zu Resultaten durch Combination und Calcül im Zimmer. Den zweiten Theil des Geschäfts habe ich bisher ganz auf mich selbst genommen, den erstern hingegen denjenigen Artillerie-Officieren übertragen, die in den Jahren 1821---1825 als Gehülfen mir zur Seite gestanden und dabei Gelegenheit gehabt hatten, nicht allein mit der Behandlung der Instrumente, sondern auch mit dem Geist und den Eigenthümlichkeiten der von mir angewandten Verfahrungsart vertraut zu werden. Diese Officiere waren: der Hauptmann \textsc{Müller}, der damalige Premier-Lieutenant \textsc{Hartmann}, und mein ältester Sohn, gegenwärtig Premier-Lieutenant im Artillerie-Regiment. Der Lieutenant \textsc{Hartmann}, welcher 1831 den Militärdienst gegen eine Anstellung bei der höhern Gewerbe-Schule in Hannover mit Hauptmanns-Charakter verliess, ist im Jahre 1834 mit Tode abgegangen.

\par\clearpage
\noindent\makebox[0pt][l]{\fontsize{13}{16}\selectfont 422}\hfill{\small\textbf{NACHLASS.}}\hfill\null\par\medskip

\begin{center}
B.\quad Bisher abgemachte Hauptdreiecksmessungen.
\end{center}

\par\noindent\hspace*{\parindent}Hätten die genannten drei Officiere ununterbrochen jedes Jahr während der ganzen tauglichen Jahreszeit sich diesem Geschäfte ausschliesslich widmen können, so würde es längst vollendet sein. Allein mancherlei Hindernisse, die theils durch die Dienstverhältnisse der Officiere, theils durch andere äussere Umstände herbeigeführt wurden, sind die Ursache gewesen, dass dieselben in einigen Jahren nur während kürzerer Zeit, in andern gar nicht daran arbeiten konnten; es kam noch dazu, ausser dem schon erwähnten Tode des Hauptmanns \textsc{Hartmann}, dass auch die fortschreitende Detailaufnahme, besonders in solchen Gegenden, wo die Kirchthürme mehr zerstreut liegen, immer noch besondere Vorbereitungsmessungen erfordert, wozu auch nur einer oder der andere jener Officiere verwendet werden konnte, dessen Zeit dann also dem Hauptgeschäfte entzogen wurde. Folgendes ist eine summarische Übersicht des in den einzelnen Jahren bisher geleisteten.

{\hangindent=.075\linewidth\hangafter=1\noindent\makebox[.075\linewidth][l]{1828}wurde die Arbeit im Spätsommer angefangen, wo Hauptmann \textsc{Müller} und Lieutenant \textsc{Gauss} die trigonometrische Messung des Eichsfeldes grösstentheils absolvirten, während der Lieutenant \textsc{Hartmann} im Amt Hunnesrück und einem Theil des Hildesheimschen Vorbereitungsmessungen für die Detailaufnahme machte.\par}

{\hangindent=.075\linewidth\hangafter=1\noindent\makebox[.075\linewidth][l]{1829.}\enspace Lieutenant \textsc{Gauss} machte zuerst einige Ergänzungsmessungen im Eichsfelde. Den grössten Theil des Jahres widmete er aber der Triangulirung in Westphalen, wo er von der Weser bis Bentheim ein Dreieckssystem ausführte. Der Hauptmann \textsc{Müller} hatte an derselben Arbeit nur eine kurze Zeit Theil nehmen können, da er bald nach dem Anfange in Folge anderer Aufträge von dem Messungsgeschäfte abberufen wurde. Die Arbeiten des Lieutenants \textsc{Hartmann} in diesem Jahre bestanden in Vorbereitungsmessungen für die Detailaufnahme, zuerst im Hildesheimschen, hernach in den Ämtern Uchte, Freudenberg und Auburg.\par}

{\hangindent=.075\linewidth\hangafter=1\noindent\makebox[.075\linewidth][l]{1830.}\enspace Trigonometrische Vermessung des östlichen Theils des Lüneburgschen durch Hauptmann \textsc{Müller} und Lieutenant \textsc{Gauss}; weitere Fortführung der im vorigen Jahre durch Westphalen gemessenen Dreieckskette bis Ostfriesland durch Lieutenant \textsc{Hartmann}.\par}

{\hangindent=.075\linewidth\hangafter=1\noindent\makebox[.075\linewidth][l]{1831}konnte nur kurze Zeit den Arbeiten gewidmet werden; Lieutenant \textsc{Gauss}\par}

\par\clearpage
\noindent\null\hfill{\small\textbf{ZUR HANNOVERSCHEN TRIANGULATION.}}\hfill\makebox[0pt][r]{\fontsize{13}{16}\selectfont 423}\par\medskip

{\leftskip=.075\linewidth\noindent vollendete die Messungen im Lüneburgschen, Lieutenant \textsc{Hartmann} die in Ostfriesland. Hauptmann \textsc{Müller} nahm gar nicht Theil.\par}

{\hangindent=.075\linewidth\hangafter=1\noindent\makebox[.075\linewidth][l]{1832}fielen die Messungen ganz aus.\par}

{\hangindent=.075\linewidth\hangafter=1\noindent\makebox[.075\linewidth][l]{1833.}\enspace Triangulirung des Landstrichs längs der Weser von der Gegend von Nienburg bis Holzminden durch Hauptmann \textsc{Müller} und Lieutenant \textsc{Gauss}; Triangulirung des Harzes durch Hauptmann \textsc{Hartmann}.\par}

{\hangindent=.075\linewidth\hangafter=1\noindent\makebox[.075\linewidth][l]{1834}konnte allein der Lieutenant \textsc{Gauss} eine auch nur sehr kurze Zeit den Geschäften widmen, theils für Vorbereitungsmessungen für die Detailaufnahme im Osnabrückschen, theils zur Ergänzung der vorigjährigen Messungen an der Weser.\par}

{\hangindent=.075\linewidth\hangafter=1\noindent\makebox[.075\linewidth][l]{1835.}\enspace Auch in diesem Jahre konnte nur Lieutenant \textsc{Gauss} auf kurze Zeit zu weitern Vorbereitungsmessungen im Osnabrückschen abkommen, diesmal unter Beihülfe des Dr. \textsc{Goldschmidt}, Observators an der Göttinger Sternwarte.\par}

{\leftskip=.075\linewidth\noindent In den beiden folgenden Jahren konnte allein der Hauptmann \textsc{Müller} sich den Messungen widmen, und zwar\par}

{\hangindent=.075\linewidth\hangafter=1\noindent\makebox[.075\linewidth][l]{1836}zur Triangulirung des Landstrichs an der Oberweser zwischen Uslar, Göttingen und Münden;\par}

{\hangindent=.075\linewidth\hangafter=1\noindent\makebox[.075\linewidth][l]{1837}zu weitern Vorbereitungen der Detailaufnahme im Osnabrückschen.\par}

{\leftskip=.075\linewidth\noindent Ausserdem machte der Hauptmann \textsc{Müller} in diesem Jahr eine vorläufige Recognoscirung der Aller-Gegend, behuf künftiger trigonometrischer Vermessung derselben.\par}

\begin{center}
C.\quad Noch fehlende Hauptdreiecksmessungen.
\end{center}

\par\noindent\hspace*{\parindent}Die bisher aufgezählten trigonometrischen Messungen hängen alle unter sich zusammen: um das \textls[120]{ganze} Königreich zu umfassen, und ein in der Hauptsache vollständiges Ganzes zu bilden, fehlt nur noch die trigonometrische Vermessung von zwei grössern Landestheilen, nemlich erstlich der Gegend rechts und links der Aller, und zweitens dem nördlichen Theile des Bremischen; in dem erstern Landestheile ist, wie schon erwähnt, eine vorläufige Recognoscirung bereits ausgeführt.

\par\noindent\hspace*{\parindent}Obgleich bei Arbeiten dieser Art unmöglich ist, den erforderlichen Zeit- und Kostenaufwand mit einiger Genauigkeit vorher zu bestimmen, da das schnellere oder langsamere Fortschreiten von so mancherlei, theils vorher nicht

\par\clearpage
\noindent\makebox[0pt][l]{\fontsize{13}{16}\selectfont 424}\hfill{\small\textbf{NACHLASS.}}\hfill\null\par\medskip

\noindent genau bekannten, theils zufälligen Umständen abhängt (z. B. besondere Terrain-Schwierigkeiten, oder Witterungszustand), so ist doch mit einiger Wahrscheinlichkeit anzunehmen, dass, wenn der Hauptmann \textsc{Müller} sich diesen Messungen allein zu widmen hat, er sie in zwei Sommern ganz oder doch grösstentheils würde vollenden können.

\par\medskip
\par\noindent\hspace*{\parindent}Die eigentlichen letzten Resultate der trigonometrischen Messungen selbst sind die \textls[120]{Zahlen}, welche die Lage der bestimmten Objecte (grösstentheils Kirchthürme) auf das schärfste festlegen, und daher einen bleibenden Werth behalten. Nach ganz vollendeter Arbeit werden dieselben geordnet, und grösserer Sicherheit wegen in mehrern Abschriften deponirt werden können, etwa eine davon in der Sternwarte. Ob solche auch durch den Druck zu publiciren sein werden, wird dann demnächst von höherer Entscheidung abhängen. Ich kann aber nicht unbemerkt lassen, dass die \textsc{Papen}sche Karte des Königreichs, welche den vollkommensten Arbeiten dieser Art in jeder Beziehung gleich kommt, und wovon 22 Blätter (also der dritte Theil des Ganzen) bereits erschienen sind, die Grundlage ihrer Genauigkeit in jenen Resultaten hat, indem dem Lieutenant \textsc{Gauss}, welcher die Graduirung und die scharfe Eintragung aller Fundamentalpositionen bei jener Karte übernommen hat, alle bisherigen Resultate zu diesem Zweck zu Gebote stehen.

\par\medskip
\par\noindent\hspace*{\parindent}Schliesslich will ich noch etwas über die Vorbereitungsmessungen zu der Detailaufnahme beifügen. Die Vollendung dieser Vorbereitungsmessungen lässt sich deswegen gar nicht im voraus bestimmen, weil sie der Natur der Sache nach nothwendig beinahe ebenso lange noch dauern müssen, wie die Detailaufnahme selbst. Denn jene bestehen eben darin, dass in Gegenden, wo die durch die Hauptdreiecksmessungen bestimmten Kirchthürme zu weit aus einander liegen und nicht zahlreich genug sind, um für jedes Messtischblatt eine hinlängliche Anzahl fester Anhaltspunkte zu liefern, besondere Signalpfähle gesetzt, und ihre Plätze noch durch kleinere Dreiecke bestimmt werden müssen. Solchen Signalpfählen kann aber kein so sicherer Schutz gegeben werden, dass auf ihr Bestehen für viele Jahre mit Gewissheit gerechnet werden könnte. Da nun aber, wenn auch nicht gerade durch das Abhandenkommen eines oder des andern einzelnen Signalpfahls, aber doch

\par\clearpage
\noindent\null\hfill{\small\textbf{ZUR HANNOVERSCHEN TRIANGULATION.}}\hfill\makebox[0pt][r]{\fontsize{13}{16}\selectfont 425}\par\medskip

\noindent durch das Verschwinden mehrerer, vor ihrer Benutzung zu der Messtischaufnahme, die ganze auf ihre Bestimmung verwandte Arbeit eine verlorne sein würde, so dürfen dergleichen Vorbereitungsoperationen immer nur höchstens ein oder ein paar Jahre früher unternommen werden, ehe die betreffende Gegend bei der Detailaufnahme an die Reihe kommt.

\par\noindent\hspace*{\parindent}Übrigens ist die Detailaufnahme selbst im Hildesheimschen, dem Eichsfelde, dem Amt Hunnesrück, den Ämtern Freudenberg, Uchte und Auburg vollendet, und im Fürstenthum Osnabrück bereits ziemlich weit vorgeschritten: allein einen genauern und vollständigen Bericht darüber wird nur der Oberst \textsc{Prott} geben können, zu dessen Ressort diese Arbeit gehört.

\par\noindent\hspace*{\parindent}Göttingen, 8. Februar 1838.

\par\medskip
\begin{center}
[Aus dem Bericht vom 5. Julius 1840 ȟber die trigonometrischen Ver-\\
messungen« im Jahre 1839.]
\end{center}

\par\noindent\hspace*{\parindent}\mbox{.\ .\ .\ .\ .\ .}\quad Bei Beurtheilung der Dreieckssysteme darf nicht übersehen werden, dass ursprünglich nicht eine allgemeine Landesvermessung beabsichtigt war, sondern zuerst nur eine Gradmessung von Göttingen bis Holstein, und sodann zunächst eine Erweiterung des Dreieckssystems zu einem Anschlusse an die \textsc{Krayenhoff}schen Dreiecke bis Ostfriesland. \textls[120]{Diesen} Zwecken gemäss waren die von mir selbst 1821---1825 gemessenen Dreiecke vom Inselsberg bis Jever \mbox{ausgewählt.\enspace\mbox{.\ .\ .\ .}}\quad Die übrigen, welche später hinzu gekommen sind, erscheinen als Abzweigungen jener Hauptdreiecke. Eine Folge dieser Entstehungsart ist, dass die Gesammtheit nicht überall in dem Maasse wie ein abgerundetes Ganzes aus Einem Guss in die Augen fällt, als der Fall gewesen sein würde, wenn eine solche Rücksicht schon von Anfang an hätte genommen werden müssen; allein der eigentliche Zweck, nemlich die scharfe Festlegung der vornehmsten sich dazu qualificirenden Punkte im ganzen Lande, ist darum nicht weniger gut \mbox{erreicht.\enspace\mbox{.\ .\ .\ .\ .\ .}}

\par\medskip
\begin{center}
[Aus einem Bericht, December 1844, »über die im Jahre 1844 ausgeführten\\
trigonometrischen Messungen.«]
\end{center}

\par\noindent\hspace*{\parindent}\mbox{.\ .\ .\ .\ .\ .}\quad Ich erlaube mir, noch einige Worte in Beziehung auf die sämmtlichen Messungen aus den vergangenen Jahren beizufügen.

\par\smallskip\noindent\hspace*{.03\linewidth}{\footnotesize IX.}\hfill{\footnotesize 54}\hspace*{.12\linewidth}\par

\par\clearpage
\noindent\makebox[0pt][l]{\fontsize{13}{16}\selectfont 426}\hfill{\small\textbf{NACHLASS.}}\hfill\null\par\medskip

\par\noindent\hspace*{\parindent}\mbox{.\ .\ .\ .\ .\ .}\quad Die Resultate [d. i. die Coordinaten] sind jedes Jahr nach Verarbeitung der Messungen in Verzeichnisse gebracht, und solcher partiellen Verzeichnisse sind sechzehn vorhanden, welche zusammen etwas über 3000 Bestimmungen enthalten, so jedoch, dass die Anzahl der Punkte selbst etwa um den siebenten Theil kleiner sein mag, indem viele Punkte, die in einem spätern Jahr nach dem Hinzukommen neuer Data schärfer oder zuverlässiger bestimmt werden konnten, in mehr als einem Verzeichnisse auftreten. Kirchthürme werden im ganzen Königreiche nicht viele ohne Bestimmung geblieben sein.

\par\noindent\hspace*{\parindent}Dass diese Verzeichnisse von allen seit 16 oder 17 Jahren vorgenommenen Detailaufnahmen sowie von den \textsc{Papen}schen Karten die Grundlage gewesen sind, braucht hier nicht weiter ausgeführt zu werden: von grosser Wichtigkeit ist aber, dass diese Zahlen, die ihren Werth behalten, so lange die Gegenstände existiren, nicht verloren gehen können. Die erwähnten Verzeichnisse werden in der Sternwarte aufbewahrt; Abschriften davon hat auch der Lieutenant \textsc{Gauss}, der alle Stammpunkte in die \textsc{Papen}schen Karten eingetragen hat, in Händen. Zu grösserer Sicherheit und bequemerm Gebrauch habe ich jetzt angefangen, die partiellen Verzeichnisse in Eins zu verschmelzen, welches demnach etwa 2600 Punkte enthalten wird.

\par\noindent\hspace*{\parindent}Späterhin könnte es vielleicht für gerathen erachtet werden, dieses Verzeichniss oder einen Auszug daraus durch den Druck zu veröffentlichen: für den Augenblick würde ich dies aber aus mehrern Gründen noch für vorzeitig halten.

\par\noindent\hspace*{\parindent}Erstlich, weil eine \textls[120]{wissenschaftlich} genügende Entwickelung von Bedeutung und allseitiger Benutzung dieser Zahlen nur nach und nach in Verbindung mit der Entwickelung der mir eigenthümlichen mathematischen Theorien gegeben werden kann, welche ich in einer Reihe einzelner Abhandlungen (etwa drei oder vier) zu liefern beabsichtige. Die erste davon ist bereits als Theil des demnächst erscheinenden Bandes der Denkschriften hiesiger Societät der Wissenschaften abgedruckt[\textsuperscript{*})] und auch einzeln in den Buchhandel gebracht; die andern werde ich nach und nach baldthunlichst nachfolgen lassen.

\par\noindent\hspace*{\parindent}Zweitens, weil die Zahlen des Verzeichnisses, obwohl hinreichend, ja überflüssig genau für jede praktische Benutzung, doch behuf der den strengsten

\par\smallskip\noindent\rule{.24\linewidth}{.4pt}\par\smallskip
{\footnotesize\noindent\hspace*{\parindent}[\textsuperscript{*}) Untersuchungen über Gegenstände der höhern Geodäsie. Erste Abhandlung. Der Königl. Societät überreicht am 23. October 1843. Band IV, S. 259/300.]\par}

\par\clearpage
\noindent\null\hfill{\small\textbf{ZUR HANNOVERSCHEN TRIANGULATION.}}\hfill\makebox[0pt][r]{\fontsize{13}{16}\selectfont 427}\par\medskip

\noindent theoretischen Forderungen entsprechenden Verschmelzung der verschiedenen Messungen in Ein System noch einige (wenn auch an sich sehr geringe) Ausfeilung und Nachhülfe zulassen.

\par\noindent\hspace*{\parindent}Drittens, weil ich die Lage der vorzüglichsten Punkte, namentlich der Kirchthürme in Städten, gern neben der Coordinatenform noch zugleich in einer andern Form, nemlich nach der geographischen Breite und Länge, beifügen möchte, welche immer einen beträchtlichen Zeitaufwand erfordernde Umformung erst nach und nach wird ausgeführt werden \mbox{können.\quad\mbox{.\ .\ .\ .\ .\ .}}

\par\medskip
\begin{center}
[3.]\\
{\large Hauptdreieckspunkte der hannoverschen Messungen.}
\end{center}

\par\smallskip
{\centering
\setlength{\tabcolsep}{0pt}
\newcommand{\GaussIXCoordinate}[7]{%
  \makebox[\linewidth][c]{%
    \makebox[1.8em][r]{\(#1\)\makebox[0pt][l]{\(\vphantom{0}#5\)}}%
    \hspace{.2em}\makebox[1.35em][r]{\(#2\)\makebox[0pt][l]{\(\vphantom{0}#6\)}}%
    \hspace{.2em}\makebox[1.55em][r]{\(#3\)\makebox[0pt][l]{\(\vphantom{0}#7\)}}%
    \( {,}#4 \)%
  }%
}
\begin{tabular}{@{\vrule width .8pt}p{.34\linewidth}|p{.23\linewidth}|p{.27\linewidth}@{\vrule width .8pt}}
\noalign{\hrule height .8pt}
\rule[-1.0\baselineskip]{0pt}{2.6\baselineskip}\makebox[\linewidth][c]{Dreieckspunkt} &
\makebox[\linewidth][c]{Breite} &
\makebox[\linewidth][c]{Länge v.\ Gött.\ Merid.} \\
\hline
\multicolumn{1}{@{\vrule width .8pt}c|}{\rule{0pt}{.45\baselineskip}} &
& \\
\hspace*{.6em}Inselsberg, hess. Dr.-P. & \GaussIXCoordinate{50}{51}{8}{619}{^{0}}{^{\prime}}{^{\prime\prime}} & \GaussIXCoordinate{-0}{31}{24}{311}{^{0}}{^{\prime}}{^{\prime\prime}} \\
\hspace*{.6em}Hohehagen & \GaussIXCoordinate{51}{28}{31}{234}{}{}{} & \GaussIXCoordinate{+0}{10}{45}{143}{}{}{} \\
\hspace*{.6em}Sternwarte in Göttingen & \GaussIXCoordinate{51}{31}{48}{028}{}{}{} & \makebox[\linewidth][c]{\(0\)} \\
\hspace*{.6em}Meridianzeichen & \GaussIXCoordinate{51}{34}{30}{309}{}{}{} & \GaussIXCoordinate{-0}{0}{0}{007}{}{}{} \\
\hspace*{.6em}Brocken & \GaussIXCoordinate{51}{48}{1}{849}{}{}{} & \GaussIXCoordinate{-0}{40}{23}{071}{}{}{} \\
\hspace*{.6em}Hils & \GaussIXCoordinate{51}{53}{53}{000}{}{}{} & \GaussIXCoordinate{+0}{6}{41}{152}{}{}{} \\
\hspace*{.6em}Lichtenberg & \GaussIXCoordinate{52}{7}{22}{020}{}{}{} & \GaussIXCoordinate{-0}{20}{33}{348}{}{}{} \\
\hspace*{.6em}Deister & \GaussIXCoordinate{52}{14}{5}{778}{}{}{} & \GaussIXCoordinate{+0}{20}{35}{702}{}{}{} \\
\hspace*{.6em}Garssen & \GaussIXCoordinate{52}{39}{40}{397}{}{}{} & \GaussIXCoordinate{-0}{12}{19}{154}{}{}{} \\
\hspace*{.6em}Scharnhorst & \GaussIXCoordinate{52}{43}{43}{000}{}{}{} & \GaussIXCoordinate{-0}{19}{1}{620}{}{}{} \\
\hspace*{.6em}Breithorn & \GaussIXCoordinate{52}{49}{45}{985}{}{}{} & \GaussIXCoordinate{-0}{19}{11}{230}{}{}{} \\
\hspace*{.6em}Falkenberg & \GaussIXCoordinate{52}{50}{52}{481}{}{}{} & \GaussIXCoordinate{+0}{4}{34}{858}{}{}{} \\
\hspace*{.6em}Hauselberg & \GaussIXCoordinate{52}{51}{36}{513}{}{}{} & \GaussIXCoordinate{-0}{14}{28}{226}{}{}{} \\
\hspace*{.6em}Elmhorst & \GaussIXCoordinate{52}{59}{42}{886}{}{}{} & \GaussIXCoordinate{+0}{18}{24}{543}{}{}{} \\
\hspace*{.6em}Steinberg & \GaussIXCoordinate{52}{59}{54}{987}{}{}{} & \GaussIXCoordinate{+0}{40}{7}{347}{}{}{} \\
\hspace*{.6em}Bottel & \GaussIXCoordinate{53}{2}{22}{919}{}{}{} & \GaussIXCoordinate{+0}{39}{22}{916}{}{}{} \\
\hspace*{.6em}Wulfsode & \GaussIXCoordinate{53}{4}{15}{113}{}{}{} & \GaussIXCoordinate{-0}{17}{48}{837}{}{}{} \\
\hspace*{.6em}Bremen, Ansgar.-Thurm & \GaussIXCoordinate{53}{4}{48}{867}{}{}{} & \GaussIXCoordinate{-1}{8}{22}{781}{}{}{} \\
\multicolumn{1}{@{\vrule width .8pt}c|}{\rule{0pt}{.45\baselineskip}} &
& \\
\end{tabular}\par
}

\par\smallskip\noindent\hfill{\footnotesize 54\textsuperscript{*}}\hspace*{.12\linewidth}\par
